%===============================================================================
% Zweck:    KTR-Seminar-Vorlage
% Erstellt: 16.10.2007
% Updated:  15.04.2013
% Autor:    U.K. / M.G.
%===============================================================================

%===============================================================================
% Zum Kompilieren pdflatex und bibtex ausführen.
%	Konfiguration: in texmaker unter Benutzer -> Benutzerbefehle einen neuen Befehl anlegen: pdflatex -synctex=1 -interaction=nonstopmode %.tex | bibtex % | makeindex %.nlo -s nomencl.ist -o %.nls | pdflatex -synctex=1 -interaction=nonstopmode %.tex | pdflatex -synctex=1 -interaction=nonstopmode %.tex
%	Entsprechende Informationen in den config/metainfo verändern
% Zur Auswahl der Sprache im folgenden Befehl
% ngerman für deutsch eintragen, english für Englisch.
%===============================================================================

% Options ngerman, english
\newcommand{\lang}{ngerman}

\documentclass[pdftex, journal, onecolumn, a4paper, 12pt, \lang]{IEEEtran}

%===============================================================================
% zentrale Layout-Angaben und Befehle
%===============================================================================
% What You should change:
% Here goes your name
\author{Author}
% and the title of your seminar
\newcommand{\subtitle}{Your Topic on the Seminar}
% the date of the submission
\date{\today}

% What is already done

\newlanguagecommand{\semester}
\addtolanguagecommand{\semester}{english}{Summer Term 2017}
\addtolanguagecommand{\semester}{ngerman}{Sommersemester 2017}

\newlanguagecommand{\ltitle}
\addtolanguagecommand{\ltitle}{english}{Fog Computing for Internet-of-Things Applications}
\addtolanguagecommand{\ltitle}{ngerman}{Fog Computing for Internet-of-Things Applications}

\title{\ltitle}
\newcommand{\supervisor}{Prof. Dr. Udo Krieger}

\gittrue
\seminartrue

%===============================================================================
% Zweck: KTR-Seminar-Vorlage in Anlehung an G. Wirtz, Lehrstuhl Praktische Informatik
%===============================================================================
%===============================================================================
% zentrale Layout-Angaben und Befehle
%===============================================================================
%
% \usepackage{german}
\usepackage{ifthen}
% Language Selection for Babel
\ifthenelse{\equal{\lang}{ngerman}}{\usepackage[german, \lang]{babel}}{\usepackage[\lang]{babel}}

\usepackage[utf8]{inputenc}
\usepackage{fancyhdr}
\usepackage[T1]{fontenc}
\usepackage{ae}
\usepackage{color}
\usepackage{amsmath}
\usepackage{amsfonts}
%

\usepackage[pdftitle={\@title},pdfauthor={\@author},pdftex,bookmarksopen,bookmarksnumbered,pdfborder=0]{hyperref}
\usepackage[pdftex]{graphicx}
\pdfcompresslevel=9

%\usepackage[dvips]{rotating}
%
% ausf\"{u}hrlichere Fehlermeldungen
\errorcontextlines=999
%
% Page-Layout: A4 aus Header
% Alternative
\setlength\headheight{14pt}
\setlength\topmargin{-15,4mm}
\setlength\oddsidemargin{-0,4mm}
\setlength\evensidemargin{-0,4mm}
\setlength\textwidth{160mm}
\setlength\textheight{252mm}
%
%% Absatzeinstellungen
\setlength\parindent{0mm}
\setlength\parskip{2ex}


\makeatletter
\renewcommand{\maketitle} {
  \begin{titlepage}
  \centering
    \begin{minipage}[t]{16cm}
      \hfill
      \begin{minipage}{12cm}
      \ifthenelse{\equal{\lang}{ngerman}}{%
        \centering
        Otto-Friedrich-Universit\"at Bamberg
        \\[12pt]%
        {\Large Professur f\"ur Informatik\\[.5em]%
        \large insbesondere Kommunikationsdienste,\\[.5em]%
        Telekommunikationsdienste und Rechnernetze}%
        }{%
        \centering
        Otto-Friedrich-University of Bamberg
        \\[12pt]
        {\Large Professorship for Computer Science,\\[.5em]%
        \large Communication Services, Telecommunication\\[.5em]%
        Systems and Computer Networks}%
        }
      \end{minipage}
      \hfill
      \begin{minipage}{3cm}
        \includegraphics[height=28mm]{config/images/logo} %height=26mm
      \end{minipage}
    \end{minipage}\\[70pt]%[50pt]
   % 
    {\Large\bf \ifthenelse{\equal{\lang}{ngerman}}{Ausarbeitung des KTR-Seminars}{Seminar on}}
    \\[36pt]
   % oder: {\LARGE Ausarbeitung des KTR-Seminars}\\[12pt]
    {\LARGE \@title}\\[80pt]
    {\Large\bf \ifthenelse{\equal{\lang}{ngerman}}{Thema:}{Topic:}}\\[36pt]
    {\LARGE\bf \subtitle}\\
    \vfill
    \begin{minipage}{\textwidth}
      \center
      \ifthenelse{\equal{\lang}{ngerman}}{Vorgelegt von:}{Submitted by:}\\
      {\Large \@author \\[18pt]}
      \ifthenelse{\equal{\lang}{ngerman}}{Betreuer:}{Supervisor:} \supervisor \\[12pt]
      Bamberg, \@date
    \end{minipage}
  \end{titlepage}
}
\makeatother
%
\usepackage[intoc]{nomencl}
\ifthenelse{\equal{\lang}{ngerman}}{%
\renewcommand{\nomname}{Abkürzungsverzeichnis}}%
{\renewcommand{\nomname}{List of Abbreviations}}
\renewcommand{\nomlabel}[1]{#1 \dotfill}
\makenomenclature

% Einbindung eines Bildes mit angegebener Breite
% #1 = label f\"{u}r \ref-Verweise
% #2 = Name des Bildes ohne Endung relativ zu Bilder-Verzeichnis
% #3 = Beschriftung
% #4 = Breite des Bildes im Dokument in cm
\newcommand{\bildw}[4]{%
  \begin{figure}[htb]%
    \centering
    \includegraphics[width=#4cm]{Bilder/#2}%
    \vskip -0.3cm%
    \caption{#3}%
    \vskip -0,2cm%
    \label{#1}%
  \end{figure}%
}
%
% Einbindung eines Bildes mit Seitenbreite
% #1 = label f\"{u}r \ref-Verweise
% #2 = Name des Bildes ohne Endung relativ zu Bilder-Verzeichnis
% #3 = Beschriftung
\newcommand{\bild}[3]{%
  \begin{figure}[htb]%
    \centering%
    \includegraphics[width=\textwidth]{Bilder/#2}%
    \vskip -0.3cm%
    \caption{#3}%
    \vskip -0,2cm%
    \label{#1}%
  \end{figure}%
}
%
\numberwithin{equation}{section}
%
%===============================================================================
% zentrale Layout-Angaben und Befehle
%===============================================================================
%

%===============================================================================
% LATEX-Dokument
%===============================================================================

\hyphenation{op-tical net-works semi-conduc-tor}
\begin{document}

\maketitle

\pagenumbering{Roman}
\setcounter{page}{2}
%
\tableofcontents
% Einstellungen f\"{u}r Literaturverzeichnis
\newpage
\addcontentsline{toc}{section}{\listfigurename}
\listoffigures
\newpage
\addcontentsline{toc}{section}{\listtablename}
\listoftables
\newpage
\printnomenclature
%===============================================================================
% LATEX-Dokument: Kapitel laden
%===============================================================================
%
\newpage
\pagenumbering{arabic}
\setcounter{page}{1}


%
% hier einzelne Kapitel mit \input{Kapitel-File} einf\"{u}gen
%
\section{Einleitung}\label{sec:ein}
Einleitung nach \autoref{sec:ein}

\section{Hauptteil}\label{sec:haupt}
\subsection{Bilder und Grafiken}\label{subsec:grafiken}
\subsubsection{Bilder}\label{subsubsec:bilder}
Bilder befinden sich im Image-Orgner und lassen sich mit \textbackslash image\{Breite\}\{Datei im Image-Verzeichnis\}\{Beschriftung\}\{Label\} einbinden. \image{3cm}{logo.png}{Uni-Logo}{img:uni} Die Referenzierung erfolg mittels \textbackslash autoref\{Label\}, also z.B. \autoref{img:uni}.
\subsubsection{Grafiken mit TikZ}
Grafiken im TikZ-Framework\footnote{\url{http://www.tn-home.de/TUGDD/Stuff/TikZ_final.pdf}} lassen sich mit dem Befehl \textbackslash scaletikzimage\{Datei im Image Verzeichnis\}\{Beschriftung\}\{Label\}\{Skalierungsfaktor\} einbinden. \scaletikzimage{tikz}{TikZ-Grafik}{img:tikz}{0.9}
\subsection{Tabellen}
Tabellen lassen sich mit dem Environment\\
\textbackslash begin\{longtable\}[H h t b c]\{Spaltendefinitionen\} ...\\
\qquad\qquad \textbackslash caption\{Tabellenunterschrift\}\\
\qquad\qquad \textbackslash label\{Label\}\\
\textbackslash end\{longtable\}\\
 definieren\footnote{\url{ftp://ftp.dante.de/pub/tex/macros/latex/required/tools/longtable.pdf}}\\
\begin{longtable}[H]{|p{0.2\textwidth}|p{0.2\textwidth}|p{0.2\textwidth}|}
\hline
A&B&C\\
\hline
\caption{Tabelle 1}
\label{tab:tab1}
\end{longtable}
\subsection{Code-Ausschnitte}
Pseudo-Code Ausschnitte lassen sich mit \textbackslash pseudo\{Name des Algorithmus\}\{Label\}\{Datei im Code-Verzeichnis\} einbinden.
\pseudo{Mittelwert}{lst:mean}{code}
\section{Zitate}
Mit \textbackslash nocite*\{\} lassen sich alle Einträge in der Bibliography ausgeben. Mit \textbackslash cite[S. xx]\{Key\} lassen sich Zitate einfügen. Z.B. \cite[S. 234]{Kurose12} \nocite*{}
\subsection{Abkürzungen}

Abkürzungen können mit \textbackslash nomenclature\{Abk\}\{\textbackslash m\{Abk\}ürzung\} \nomenclature{Abk}{\m{Abk}ürzung} angegeben werden. Diese werden alphabetisch sortiert in ein Abkürzungsverzeichnis aufgenommen.


%
%===============================================================================
% LATEX-Dokument: Literaturverzeichnis
%===============================================================================
%
\newpage
\phantomsection
% Einstellungen f\"{u}r Literaturverzeichnis
\addcontentsline{toc}{section}{\bibname}

\bibliographystyle{IEEEtran}
% argument is your BibTeX string definitions and bibliography database(s)
\bibliography{literature/bib}
% Nutzung von Bibtex:
% hier den bib-file einbinden
%
% GATHER{bibfile.bib}
% \footnotesize
% \bibliography{bibfile}
% ansonsten: bbl als tex Datei einbinden
 %\input{KTR-Seminar-Literatur.tex}
%===============================================================================
% LATEX-Dokument: Literaturverzeichnis
%===============================================================================
%
\end{document}
